\section{Introduction}
A composition is a nonnegative vector summing to one. An uncertainty region around a compositional prediction is useful only if it covers the intended target. A classifier's probability vector can be used to predict a \emph{class label}; an observed vector of proportions can instead be the response; a latent conditional probability vector is a third target. Covering one does not establish coverage of the others \citep{angelopoulos2020raps,javanmardi2024credal}.

We study joint coverage of a future \emph{observed composition}, using split-conformal scores with separately fitted predictors. Exchangeable calibration and test scores give a finite-rank marginal guarantee \citep{papadopoulos2002inductive,lei2016distribution,angelopoulos2021gentle}. This leaves practical choices of geometry, adaptation and calibration granularity. Equal marginal guarantees need not allocate uncertainty similarly across the simplex or across observable groups. Raw thresholds in different geometries cannot determine which region is smaller.

Our questions are: (Q1) how do geometry and training-only scaling affect joint reliability and original-simplex efficiency under exchangeability? (Q2) how much minority undercoverage can marginal validity hide, and what do scaling and group calibration cost? (Q3) how do boundary zeros, calibration scarcity and a specified calibration--test mismatch delimit applicability? Synthetic observations make the whole-vector event observable and let mechanisms vary separately from calibration. All methods share partitions and a fitted center, isolating score/calibration choices rather than comparing score-specific optimal predictors.

Before confirmation, we hypothesized that joint scores attain their finite-rank marginal target while naive coordinate calibration need not; area orderings depend on the mechanism; and scaling or group calibration can reduce hidden minority failure at an efficiency or sample-size cost. Zeros and shift delimit domain and guarantee claims. We executed 19 frozen cells with 100 independently regenerated/refitted trials each, separately auditing numerical area uncertainty. The contribution is a reproducible evaluation with controlled positive, negative and unresolved findings, not a new conformal algorithm. The results show a geometry-ordering reversal, severe hidden minority failure, conservative full-domain outcomes with scarce group calibration, and the inability of full-dimensional area to diagnose zero-face coverage.
